% =====================================================================
%  Memory capacity as a retrieval limit
%  Muhammad Rakibul Islam -- Independent Research
%
%  Zenodo DOI: 10.5281/zenodo.21881612
%  Code:       https://github.com/rakib-nyc/axon
%
%  Styling is original. No journal class or template is used.
%
%  ONE ITEM requires author input before deposit: search % VERIFY
% =====================================================================

\documentclass[10pt,twocolumn,a4paper]{article}

\usepackage[T1]{fontenc}
\usepackage[utf8]{inputenc}
\usepackage[english]{babel}
\usepackage{mathptmx}
\usepackage[scaled=0.90]{helvet}
\usepackage{courier}

\usepackage[a4paper,left=15mm,right=15mm,top=18mm,bottom=22mm,columnsep=6mm]{geometry}

\usepackage{amsmath,amssymb}
\usepackage{booktabs}
\usepackage{array}
\usepackage{microtype}
\usepackage{caption}
\usepackage{enumitem}
\usepackage{stfloats}
\usepackage{titlesec}
\usepackage{fancyhdr}
\usepackage{xcolor}
\usepackage{xurl}
\usepackage[hidelinks,breaklinks]{hyperref}

\definecolor{accent}{RGB}{17,84,102}
\definecolor{rule}{RGB}{17,84,102}

\hypersetup{
  pdftitle={Memory capacity as a retrieval limit: direct synaptic measurement across a capacity range in a sparse binary assembly network},
  pdfauthor={Muhammad Rakibul Islam},
  pdfsubject={Associative memory, storage capacity, retrieval dynamics},
  pdfkeywords={associative memory; assembly codes; storage capacity; retrieval dynamics; bounded synapses; sparse binary networks}
}

% --- headings ---
\renewcommand{\thesection}{\Roman{section}}
\renewcommand{\thesubsection}{\thesection-\Alph{subsection}}
\titleformat{\section}[block]
  {\normalsize\bfseries\sffamily\color{accent}}{\thesection.}{0.6em}{\MakeUppercase}
\titlespacing*{\section}{0pt}{2.4ex plus 0.6ex minus 0.2ex}{1.2ex plus 0.2ex}
\titleformat{\subsection}[block]
  {\normalsize\bfseries\sffamily}{\Alph{subsection}.}{0.5em}{}
\titlespacing*{\subsection}{0pt}{1.9ex plus 0.5ex minus 0.2ex}{0.9ex plus 0.2ex}

\captionsetup[table]{font=footnotesize,labelfont={bf,sf,color=accent},skip=5pt,
                     justification=justified,labelsep=period}
\setlist[enumerate]{leftmargin=*,itemsep=2pt,topsep=3pt}
\setlist[itemize]{leftmargin=*,itemsep=2pt,topsep=3pt}

\newcommand{\pct}{\%}
\newcommand{\doilink}{10.5281/zenodo.21881612}

\setlength{\dbltextfloatsep}{14pt plus 2pt minus 2pt}
\renewcommand{\dbltopfraction}{0.9}
\renewcommand{\dblfloatpagefraction}{0.7}
\setcounter{dbltopnumber}{2}

\fancypagestyle{plain}{%
  \fancyhf{}\renewcommand{\headrulewidth}{0pt}%
  \fancyfoot[C]{\footnotesize\color{accent}\thepage}%
}
\pagestyle{plain}

% --- lead-in block macro (abstract / index terms) ---
\newcommand{\leadin}[1]{\noindent{\sffamily\bfseries\color{accent}\footnotesize
  \MakeUppercase{#1}}\hspace{0.8em}}

\begin{document}

% =====================================================================
\twocolumn[
  \begin{@twocolumnfalse}
  \vspace*{-8mm}
  \noindent{\color{rule}\rule{\textwidth}{1.2pt}}

  \vspace{3mm}
  \noindent{\sffamily\bfseries\LARGE\color{accent}
   Memory Capacity as a Retrieval Limit:\\[1mm]
   Direct Synaptic Measurement Across a Capacity Range\\[1mm]
   in a Sparse Binary Assembly Network\par}
  \vspace{5mm}
  \noindent{\sffamily\normalsize MUHAMMAD RAKIBUL ISLAM\par}
  \vspace{1mm}
  \noindent{\footnotesize Independent Research\quad$\cdot$\quad\texttt{imrakibul@gmail.com}\par}
  \vspace{2mm}
  \noindent{\footnotesize
   Preprint, August 2026\quad$\cdot$\quad
   DOI \href{https://doi.org/\doilink}{\doilink}\quad$\cdot$\quad
   Code \url{https://github.com/rakib-nyc/axon}\par}
  \vspace{1mm}
  \noindent{\footnotesize Released under a Creative Commons Attribution 4.0 International
   (CC BY 4.0) licence.\par}
  \vspace{3mm}
  \noindent{\color{rule}\rule{\textwidth}{0.5pt}}

  \vspace{5mm}
  \end{@twocolumnfalse}
]

% =====================================================================
\leadin{Abstract}
Storage capacity in associative memory is normally treated as a property of the
synaptic substrate: patterns superpose until interference makes them unrecoverable. This
paper reports a sparse binary assembly network in which that account does not hold over
the range where recall fails, established by measuring the synapses directly rather than
inferring their state from behaviour. The network has $131{,}072$ neurons and
$1.07\times10^{9}$ binary synapses with bounded $4$-bit counters, contrastive learning
and an exact $k$-winners-take-all readout, and runs on a single consumer GPU. Over the
range $24$ to $60$ stored patterns, at five random seeds per load, the sustained-recall
ratio falls by a factor of $3.7$ while the functional density of recurrent synapses
within stored assemblies remains constant at $80.7$--$82.2\pct$. The constancy holds
separately for neurons exclusive to one assembly ($82.7\pct \to 81.2\pct$) and for
neurons shared across several ($69.3\pct \to 82.1\pct$, rising) while the shared fraction
of each assembly grows from $4.8\pct$ to $64.4\pct$, so it is not an artifact of
averaging over a degrading sub-population. Retrieval initiation is likewise
load-invariant in two independent measurements, and at intermediate load recurrent drive
is maintained at unity gain while the retrieved population nonetheless drifts away from
the stored assembly. The failure mode is identity drift into a structured neighbourhood:
neurons visited during free running but outside the target assembly receive
within-assembly recurrent input at $6.9\times$ the population baseline at low load,
falling to $1.8\times$ at high load as the composition of that neighbourhood shifts from
unallocated neurons to neurons belonging to competing assemblies. Because recall depends
on how long retrieval is allowed to run, capacity is reported as a pair: sustained
capacity $\approx 36$ patterns and transient capacity $\approx 200$.

\vspace{3mm}
\leadin{Index Terms}
Assembly codes, associative memory, bounded synapses, retrieval dynamics, sparse binary
networks, storage capacity.

\vspace{4mm}

% =====================================================================
\section{Introduction}

Storage capacity in associative memory is classically an interference problem. In
Hopfield and Willshaw-type models, each stored pattern contributes to a shared synaptic
matrix, and the cross-talk between patterns grows until retrieval fails
\cite{knoblauch2010}. In models with bounded synapses the same logic is sharper still:
because synaptic efficacy cannot grow without limit, storing new memories necessarily
degrades old ones, and memory lifetime becomes the central quantity
\cite{fusiabbott2007}. In both traditions the substrate is where capacity lives.

Recent theoretical work has questioned whether this is the whole story. Tamamori's
analysis of kernel Hopfield networks concludes that the practical storage limit is
governed not by a lack of geometric separability in feature space but by the loss of
dynamical stability against cross-talk noise \cite{tamamori2026,tamamori2025}, and Clark
\cite{clark2026} shows explicitly that transient recall and persistent (fixed-point)
recall are distinct capacities, with patterns retrievable above capacity despite the
absence of stable attractors.

Those analyses proceed from dynamics and theory. They do not open the synaptic matrix and
inspect it while behaviour fails. The system studied here permits exactly that: the
synapses are bounded $4$-bit counters whose functional state is a single bit, stored
assemblies are explicit sets of neurons, and both can be enumerated exactly at any point
during or after learning. We can therefore ask a question that is usually inaccessible:
\emph{when recall degrades with load, has the substrate that supports it degraded?}

The answer in this system is no, over the range in which the behavioural transition
occurs. Section~\ref{sec:measured} presents that measurement and its controls,
Section~\ref{sec:ruledout} the alternative explanations we tested and excluded, and
Section~\ref{sec:discussion} what the result does and does not support.

% =====================================================================
\section{System Under Study}
\label{sec:system}

Full implementation detail is in Section~\ref{sec:methods}. This section gives only what
is needed to read the results.

\subsection{Population and connectivity}
The network comprises $4096$ columns of $32$ neurons, $N = 131{,}072$ neurons in total.
Columns $0$--$511$ ($16{,}384$ neurons) form a clamped input area; the remaining $3584$
columns ($114{,}688$ neurons) form the recurrent pool in which assemblies are stored.
Each neuron has $32$ dendritic branches, each sampling a fixed group of $256$ presynaptic
neurons, giving a fan-in of $8192$ and $1.07\times10^{9}$ synapses in total. Four branches
sample the input area (the afferent path); $28$ sample the recurrent pool.

\subsection{Synapses and dynamics}
Each synapse is a bounded $4$-bit counter stored as four bit-planes. The \emph{functional}
weight seen by propagation is the most significant bit: a synapse either transmits or does
not. Learning updates the counter; only threshold crossings are visible to the dynamics.
Under the contrastive rule the counters are clipped to a symmetric band about the
threshold. Total weight memory is $512$\,MB.

At each timestep each neuron accumulates $\mathrm{popcount}(w \wedge s)$ over its
branches, where $w$ is its functional weight vector and $s$ the presynaptic spike pattern.
A global threshold selects the $64$ columns with the highest drive; within each selected
column an exact $k$-winners-take-all with $k=24$ selects the firing neurons. This
selection rule and the assembly-formation scheme it supports are closest in spirit to the
assembly calculus of Papadimitriou et al.\ \cite{papadimitriou2020}, which combines a
$k$-cap operation with Hebbian plasticity on random connectivity. The active set is
therefore pinned at $64\times24 = 1536$ neurons by construction --- a fact that matters
when interpreting turnover (Section~\ref{sec:drift}).

\subsection{Learning}
Learning is contrastive: with the cue clamped, co-active pairs are potentiated; with the
cue released, co-active pairs are depressed. An intrinsic excitability term suppresses
recently active neurons during learning and is disabled at readout. Unless stated
otherwise all results use the configuration given in Section~\ref{sec:methods}, trained
for $40$ epochs.

\subsection{Metrics}
\emph{Overlap} between neuron sets $A$ and $B$ is $|A\cap B|/\sqrt{|A|\,|B|}$.
\emph{Recall} is the overlap between the retrieved state and the assembly stored for that
cue, measured at a $75\pct$ cue. \emph{Best-other} is the maximum overlap with any
non-target assembly. \emph{Margin} is recall minus best-other. \emph{Sustained ratio} is
the overlap with the stored assembly at the last free step divided by that at the first
free step, over a $30$-step free run; the sustained criterion is a ratio of at least
$0.5$. Every figure reported below is accompanied by an untrained control run under
identical settings.

The sustained ratio is a quotient and becomes uninterpretable when its denominator
collapses, so it is reported only under three conditions: (a) mean free-phase recurrent
drive to assembly members exceeds twice the untrained floor of $29$, (b) recall exceeds
best-other, and (c) the matched untrained control does not itself satisfy the sustained
criterion. Each condition excludes a distinct degenerate case: (a) collapse of network
activity, (b) stable retrieval of a state other than the target, and (c) mechanisms that
hold any state regardless of what is stored.

% =====================================================================
\section{Results}
\label{sec:measured}

\subsection{Recurrent synaptic density is constant across the behavioural transition}
\label{sec:density}

We measure \emph{within-assembly functional density}: among all synapses whose
presynaptic and postsynaptic neurons both belong to a given stored assembly, the fraction
whose counter is above threshold. This is computed exactly, over the $28$ recurrent
branches only, for every assembly.

Table~\ref{tab:capacity} gives the full capacity sweep and Table~\ref{tab:seeds} the
five-seed replication over the range in which the behavioural transition occurs.

\begin{table*}[!t]
\centering
\caption{Capacity sweep, one run per load. ``within'' is within-assembly functional
density; ``drive'' is mean recurrent drive to assembly members at the first free step;
``criterion'' records whether the sustained ratio satisfies the three conditions of
Section~\ref{sec:system}. Untrained controls were run at every load and are flat:
within-assembly density $6.3613\pct$ at all loads, margin $-0.029$ to $-0.079$, drive
$28$, none satisfying the criterion. Single-run sustained ratios should be read against
the replicated values in Table~\ref{tab:seeds}.}
\label{tab:capacity}
\small
\begin{tabular}{rrrrrrrl}
\toprule
patterns & within (\pct) & recall & best-other & margin & sustained & drive & criterion \\
\midrule
12  & 82.09 & 0.746 & 0.042 & $+0.704$ & 0.661 & 118 & met \\
24  & 82.41 & 0.803 & 0.089 & $+0.714$ & 0.597 & 119 & met \\
36  & 81.34 & 0.802 & 0.129 & $+0.673$ & 0.652 & 120 & met \\
48  & 82.07 & 0.705 & 0.167 & $+0.537$ & 0.284 & 118 & met \\
60  & 81.78 & 0.674 & 0.182 & $+0.492$ & 0.207 & 116 & met \\
74  & 78.40 & 0.551 & 0.195 & $+0.356$ & 0.090 & 111 & met \\
90  & 74.27 & 0.421 & 0.182 & $+0.239$ & 0.046 & 104 & met \\
140 & 64.94 & 0.347 & 0.178 & $+0.169$ & 0.027 & 93  & met \\
200 & 63.43 & 0.237 & 0.186 & $+0.051$ & 0.020 & 96  & met \\
300 & 64.52 & 0.137 & 0.179 & $-0.041$ & 0.024 & 102 & \emph{not met} \\
\bottomrule
\end{tabular}
\end{table*}

\begin{table*}[!t]
\centering
\caption{Five-seed replication. Seeds vary the random weight initialisation, the
connectivity and the stimulus set jointly. Density is highly reproducible; the sustained
ratio is not. CV is the coefficient of variation of the sustained ratio.}
\label{tab:seeds}
\small
\begin{tabular}{rlllr}
\toprule
patterns & sustained ratio (mean $\pm$ SD) & margin & within-assembly density (\pct) & sustained CV \\
\midrule
24 & $0.664 \pm 0.051$ & $+0.755 \pm 0.026$ & $82.23 \pm 0.17$ & $7.7\pct$ \\
36 & $0.523 \pm 0.096$ & $+0.681 \pm 0.012$ & $81.42 \pm 0.42$ & $18.4\pct$ \\
42 & $0.422 \pm 0.066$ & $+0.620 \pm 0.026$ & $80.69 \pm 0.63$ & $15.6\pct$ \\
48 & $0.326 \pm 0.083$ & $+0.577 \pm 0.033$ & $81.15 \pm 0.73$ & $25.5\pct$ \\
60 & $0.178 \pm 0.111$ & $+0.499 \pm 0.007$ & $81.49 \pm 1.12$ & $62.4\pct$ \\
\bottomrule
\end{tabular}
\end{table*}

Across patterns $24$ to $60$ the sustained ratio falls by a factor of $3.7$ while
within-assembly density stays within the band $80.7$--$82.2\pct$, with per-load standard
deviations no larger than $1.12$ percentage points. The corresponding five-seed sustained
ratio at $12$ patterns is $0.672$. Density does eventually erode --- $78.40\pct$ at $74$
patterns, $74.27\pct$ at $90$, $64.94\pct$ at $140$ --- but that erosion begins well past
the range in which the behavioural transition occurs, and is therefore a consequence of
continued loading rather than its cause.

\subsection{Capacity is reported as a pair}
Recall depends on how long free running is allowed to continue.
\textbf{Sustained capacity} --- the highest load at which the mean trace still retains at
least half its initial free-run overlap after $30$ free steps --- is $\approx 36$
patterns; the next load tested, $42$, does not satisfy the criterion
($0.422 \pm 0.066$). It is a soft crossing rather than a boundary: the $\pm 1$\,SD
intervals of adjacent loads overlap throughout ($36$: $[0.427,0.619]$; $42$:
$[0.356,0.488]$; $48$: $[0.243,0.409]$), so no single run identifies which side of the
criterion a network falls on. \textbf{Transient capacity} --- the highest load at which
margin at a $75\pct$ cue remains positive and the criterion is met --- is $\approx 200$
patterns. The gap between these two numbers is a property of the measurement horizon, not
a discrepancy.

\subsection{The constancy is not an averaging artifact}
\label{sec:subgroups}

A mean can conceal a degrading sub-population. The natural candidate is neurons belonging
to more than one assembly, which are the most exposed to competing plasticity. We
therefore split each assembly's members into \emph{shared} (belonging to at least one
other assembly) and \emph{exclusive}, and measured density separately
(Table~\ref{tab:subgroups}).

\begin{table*}[!t]
\centering
\caption{Within-assembly density split by whether the postsynaptic member is shared with
another assembly. Both groups hold; the shared group improves.}
\label{tab:subgroups}
\small
\begin{tabular}{rrrrr}
\toprule
patterns & shared fraction & density, shared (\pct) & density, exclusive (\pct) & difference \\
\midrule
12 & 0.048 & 69.26 & 82.69 & $-13.4$ \\
24 & 0.172 & 78.46 & 83.16 & $-4.7$ \\
36 & 0.366 & 80.39 & 81.82 & $-1.4$ \\
48 & 0.544 & 82.78 & 81.31 & $+1.5$ \\
60 & 0.644 & 82.09 & 81.24 & $+0.8$ \\
\midrule
48 (untrained) & 0.496 & 6.38 & 6.35 & $+0.03$ \\
\bottomrule
\end{tabular}
\end{table*}

The shared fraction of each assembly rises from $4.8\pct$ to $64.4\pct$ across the load
range. Exclusive-member density is flat ($82.69 \to 81.24\pct$) and shared-member density
\emph{rises} ($69.26 \to 82.09\pct$). At low load the few shared neurons are atypical
stragglers; by high load sharing is the norm and their density matches the rest. There is
no degrading sub-population hidden inside the mean.

\subsection{Retrieval initiation is load-invariant}
\label{sec:initiation}

Two independent quantities measured at the moment of cue release do not vary with load.

First, the overlap lost when the cue is removed. Writing $o_5$ for the overlap at the last
clamp-influenced step and $o_6$ for the first genuinely free step, the ratio $o_6/o_5$ is
$0.706 \pm 0.002$ at $12$ patterns, $0.728 \pm 0.012$ at $36$, and $0.694 \pm 0.022$ at
$60$.
% VERIFY: insert the number of seeds behind these three standard deviations.
Releasing the cue costs a fixed $\approx 29\pct$ of assembly overlap regardless of how
much is stored.

Second, the recurrent drive received by assembly members at the first free step is
$116$--$120$ over patterns $12$ to $60$ (Table~\ref{tab:capacity}); restricted to the
exclusive-member subgroup the same quantity is $119$--$126$. We scope this claim to the
range of the behavioural transition: beyond it the initial drive does decline ($93$ at
$140$ patterns).

Over that range, retrieval initiation is load-blind; load acts only on what follows.

\subsection{At intermediate load, drive is maintained while the trace decays}
\label{sec:gain}

We define \emph{loop gain} as the ratio of recurrent drive to assembly members between
steps of the same parity, two steps apart. Same-parity comparison is required because the
free-running dynamics contain a period-2 component (Section~\ref{sec:period2}); a
step-to-step ratio would conflate the oscillation with the trend.

\begin{table*}[!t]
\centering
\caption{Loop gain per two-step period across the free run, five seeds. Gain relaxes
towards unity at every load, approaching from above at low load and from below at high
load. Because gain is not constant across the run, an early-window estimate cannot be
extrapolated over the full $30$ steps.}
\label{tab:gaintraj}
\small
\begin{tabular}{rrrrrrr}
\toprule
patterns & $t_{6\to8}$ & $t_{10\to12}$ & $t_{14\to16}$ & $t_{18\to20}$ & $t_{22\to24}$ & $t_{26\to28}$ \\
\midrule
12 & 1.074 & 1.035 & 1.014 & 0.988 & 1.006 & 0.986 \\
36 & 0.998 & 1.016 & 0.996 & 1.008 & 0.977 & 0.998 \\
60 & 0.892 & 0.950 & 0.910 & 0.940 & 0.965 & 0.971 \\
\bottomrule
\end{tabular}
\end{table*}

The row for $36$ patterns is the important one. Loop gain is $0.977$--$1.016$ --- unity
within noise --- across the entire free run, and the active population is pinned at $1536$
neurons by the readout. Yet overlap with the stored assembly decays from $0.720$ to
$0.520$ over the same interval. At intermediate load the trace degrades without any loss
of recurrent drive. Whatever is failing, it is not amplitude.

\subsection{The failure mode is identity drift into a structured neighbourhood}
\label{sec:drift}

Since population size is fixed by the readout and drive is maintained, we measured which
neurons are active and how that set changes.

\paragraph{Turnover and leak.}
Per free step, the fraction of the active set replaced is $31.3\pct$ at $12$ patterns,
$32.9\pct$ at $36$, and $51.6\pct$ at $60$; the untrained control replaces $88$--$96\pct$
per step. Over the same interval the overlap with the stored assembly is lost at
$0.149 \pm 0.062$, $1.199 \pm 0.639$ and $5.162 \pm 1.844$ percent per step. The ratio of
churn to leak is $255 \pm 148$, $37 \pm 24$ and $11 \pm 4$; these are reported as an
ordering rather than as point values, the low-load figure being poorly resolved. What is
well resolved is that turnover is essentially identical at $12$ and $36$ patterns
($31.3\pct$ versus $32.9\pct$) while leak differs eightfold. Load changes the directedness
of churn, not its magnitude.

\paragraph{The active set is not a fixed population.}
The union of all neurons active at any point during the $30$ free steps is $3265 \pm 60$
at $12$ patterns, $4701 \pm 702$ at $36$, and $11{,}257 \pm 1410$ at $60$, against an
active set of $1536$ --- ratios of $2.13$, $3.06$ and $7.33$. The untrained control visits
$30{,}199 \pm 1611$ ($19.66\times$). The trace therefore cycles through a population
substantially larger than the readout displays at any instant, and the learned constraint
bounds the size of that population --- a ninefold tightening relative to untrained ---
rather than fixing its membership. The union co-varies monotonically with load, as does
the sustained ratio; we report it as a descriptive quantity and do not treat the
correlation between them as evidence of mechanism, since any two monotone functions of
load must correlate.

\paragraph{The visited population is structured.}
Table~\ref{tab:union} characterises it. Neurons visited during free running but outside
the stored assembly receive within-assembly recurrent input at $59.68\pct$ at $12$
patterns, against a population baseline of $8.62\pct$: an enrichment of $6.9\times$. In
the untrained control the same two numbers are $6.28\pct$ and $6.25\pct$, an enrichment of
$1.00\times$, confirming that the measurement is not an artifact of the procedure.

\begin{table*}[!t]
\centering
\caption{Characterisation of the population visited during free running. ``dens.\ out'' is
the functional density of recurrent synapses from stored assembly members onto
visited-but-non-member neurons; ``dens.\ all'' is the same quantity for the whole
recurrent pool. ``mult.'' is the mean number of assemblies a neuron belongs to, given for
the outside group and for the whole pool.}
\label{tab:union}
\footnotesize
\setlength{\tabcolsep}{4.5pt}
\begin{tabular}{lrrrrrrr}
\toprule
 & in-assy. & outside & \pct\ outside & dens.\ in (\pct) & dens.\ out (\pct) & dens.\ all (\pct) & mult.\ out / pool \\
\midrule
12 patterns & 1531 & 1733 & 53.1 & 82.13 & 59.68 & 8.62 & $0.12 / 0.16$ \\
36 patterns & 1536 & 3166 & 67.3 & 81.43 & 25.61 & 8.74 & $1.00 / 0.48$ \\
60 patterns & 1536 & 9721 & 86.4 & 81.49 & 16.70 & 9.10 & $1.53 / 0.80$ \\
\midrule
36, untrained & 613 & 29{,}585 & 98.0 & 6.38 & 6.28 & 6.25 & $0.56 / 0.48$ \\
\bottomrule
\end{tabular}
\end{table*}

Two measurements point the same way from opposite directions. The enrichment of
within-assembly input onto visited non-members falls
$6.9\times \to 2.9\times \to 1.8\times$ with load. Simultaneously, the
assembly-multiplicity of those neurons inverts: at $12$ patterns they are \emph{less}
likely than average to belong to another assembly ($0.12$ versus a population mean of
$0.16$), that is, they are mostly unallocated neurons that happen to be well wired to this
assembly; by $60$ patterns they are nearly twice as likely ($1.53$ versus $0.80$), that
is, they belong to competing assemblies.

Load does not merely widen the population the trace explores; it changes that
population's composition from assembly-adjacent to inter-assembly.

\subsection{Allocation}
\label{sec:allocation}

Assemblies are allocated with very low mutual overlap. At $12$ patterns the union of all
assigned neurons is $18{,}015$ against a sum of assembly sizes of $18{,}456$ (excess
$441$, i.e.\ $2.4\pct$ reuse), and the mean pairwise assembly Jaccard index is $0.00222$
against a chance value of $0.00674$ for independently drawn sets of the same size ---
$3.0\times$ below chance. The untrained control gives $0.00946$, which is $1.4\times$
\emph{above} chance.

The stricter test accounts for forced sharing. At high load, storing many assemblies of
$1536$ neurons in a pool that the allocator uses only partially requires reuse, so the
appropriate null is $m^{2}/U$, where $m$ is assembly size and $U$ the realised union. At
$140$ patterns the measured mean pairwise Jaccard index is $0.00647$ against a
sharing-adjusted null of $0.00820$ ($0.79\times$); at $300$ patterns it is $0.00777$
against $0.00747$ ($1.04\times$), at a mean neuron multiplicity of $4.45$. Allocation
overlap therefore remains at or below what forced sharing requires even in the heavy-reuse
regime.

Competitive learning and $k$-winners-take-all are standard mechanisms for producing
decorrelated sparse codes \cite{makhzani2015,sacouto2023}, so sub-chance overlap is close
to expected; what is reported here is that decorrelation persists into a regime where
sharing is geometrically unavoidable.

\paragraph{Residual best-other overlap is a tail statistic.}
At $140$ and $300$ patterns, best-other is $0.178$ and $0.179$, but the \emph{mean}
overlap with all non-target assemblies is $0.025$ and $0.026$ --- close to the
sharing-adjusted null, and equal to the untrained floor of $0.025$. Best-other is
therefore driven by a small number of high-overlap pairs rather than by uniform confusion.
Those pairs correlate with shared cue columns (Pearson $r = +0.403$ and $+0.355$ trained,
versus $+0.098$ and $+0.097$ untrained), so the network learns a mapping from input
geometry into assembly geometry. Retrieval failure, however, does not track that mapping
(Section~\ref{sec:ruledout}).

\subsection{Period-2 alternation}
\label{sec:period2}

Free-running states overlap more with the state two steps earlier than with the
immediately preceding state, in trained networks only: at $12$ patterns $J(t,t{-}2)$ runs
from $0.474$ to $0.602$ across the free run against $J(t,t{-}1)$ of $0.485$ to $0.526$,
whereas the untrained control gives $0.025$--$0.034$ and $0.021$--$0.036$ --- equal within
noise. The alternation is therefore a learned dynamical property rather than a sampling
artifact.

It is not, however, the mechanism behind the wide visited population. If free running were
a clean two-cycle over disjoint sets, each parity stream would visit $\approx 1536$
neurons. Measured, each parity stream alone visits $1.90$--$1.93\times$ the active set at
$12$ patterns, and the two parity sets overlap at $0.892 \pm 0.008$: they are nearly the
same population. Parity overlap declines monotonically with load
($0.892 \to 0.814 \to 0.746 \to 0.672 \to 0.617$ over patterns $12$ to $60$; untrained
$0.299$), so alternation strengthens as the trace degrades. It accompanies the failure
rather than causing it.

% =====================================================================
\section{Alternative Explanations Tested and Excluded}
\label{sec:ruledout}

Each of the following is a candidate explanation for the load-dependent decline in recall.
Each was tested and excluded. Untrained controls accompany every entry.

\paragraph{Hub capture.}
Hypothesis: a small set of highly connected neurons wins the competition regardless of
cue. Excluded. The functional in-degree distribution has coefficient of variation $0.150$
and $p_{99}/p_{50} = 1.25$ --- no heavy tail. Neurons winning in at least $11$ of $12$
free runs have \emph{lower} in-degree ($499.2$) than neurons that never win ($680.8$), a
ratio of $0.733$, inverted relative to the prediction.

\paragraph{Synaptic erosion at the boundary.}
Hypothesis: contrastive depression driven by other patterns erodes the recurrent scaffold.
Excluded over the relevant range: within-assembly density is $82.09$, $82.41$, $81.34$ and
$82.07\pct$ at $12$, $24$, $36$ and $48$ patterns while the five-seed sustained ratio falls
from $0.672$ to $0.326$ (Section~\ref{sec:density}). Erosion occurs later and is a
consequence, not the cause.

\paragraph{Allocation interference.}
Excluded: measured pairwise assembly overlap is $0.79$--$1.04\times$ the sharing-adjusted
null at up to $4.45\times$ forced multiplicity (Section~\ref{sec:allocation}).

\paragraph{Input-space collisions.}
Hypothesis: cues drawn from a limited input space overlap, so assemblies must overlap.
Excluded by a controlled comparison in which pool size and afferent drive were held fixed
and only the expected pairwise cue overlap was varied fourfold (expected shared cue
columns $2.0$ versus $8.0$): margin $-0.080$ versus $-0.086$, mean non-target overlap
$0.039$ versus $0.041$. Indistinguishable.

\paragraph{Pool size.}
Hypothesis: capacity is limited by the number of available neurons. Excluded at fixed code
sparsity: over a sixfold pool range ($32{,}768$ to $200{,}704$ recurrent neurons, with the
number of active columns scaled to hold the active-column fraction constant at $1.79\pct$),
recall at $140$ patterns is $0.300$, $0.345$, $0.347$ and $0.325$ --- flat --- while
within-assembly density is likewise flat at $60.7$--$65.3\pct$.

\paragraph{Averaging artifact.}
Excluded by the shared/exclusive split of Section~\ref{sec:subgroups}.

\paragraph{Drive magnitude at intermediate load.}
Excluded: loop gain is $0.977$--$1.016$ across the entire free run at $36$ patterns while
overlap decays $0.720 \to 0.520$ (Section~\ref{sec:gain}).

\paragraph{Readout-side rescue.}
Two interventions on frozen weights were tested and neither restores sustained activity in
a way attributable to the stored memory. Reducing $k$ does not restore it (best sustained
ratio $0.334$ at $k=20$, worse below $k=16$); it lowers total drive rather than sharpening
competition, and the relevant competition is between columns while $k$ operates within one.
Adding a facilitation term favouring recently active neurons does raise the sustained ratio
($0.284 \to 0.814$ at $48$ patterns), but the matched untrained control also sustains, at
$1.000$: the mechanism holds any state. Measured as learned contribution (trained margin
minus untrained margin) it \emph{reduces} the learned signal at low load, from $+0.733$ to
$+0.268$ at $12$ patterns, while raising it at high load, from $+0.038$ to $+0.124$ at
$300$. It is a load-dependent trade rather than a fix.

% =====================================================================
\section{Discussion}
\label{sec:discussion}

The central result is narrow: in this system, over the load range in which sustained recall
collapses, the synapses that support recall do not change. The supporting evidence is that
the constancy survives disaggregation into the two sub-populations where degradation would
be expected to appear first, that two independent measurements of retrieval initiation are
also load-invariant, and that at intermediate load recurrent drive is maintained at unity
gain while the trace decays anyway.

What replaces substrate degradation as the explanation is identity drift, and the
characterisation in Section~\ref{sec:drift} makes that specific rather than merely
negative: the trace visits a population two to seven times the size of the instantaneous
active set, that population is strongly enriched for within-assembly input at low load, and
load shifts its composition from unallocated assembly-adjacent neurons towards neurons
belonging to competing assemblies. The composition shift, supported by two measurements
pointing in opposite directions --- connectivity enrichment falling, assembly multiplicity
rising --- is the most specific mechanistic statement the data support.

\subsection{Relation to bounded-synapse theory}

The natural comparison is to the bounded-synapse literature. Fusi and Abbott
\cite{fusiabbott2007} show that with bounded synapses, storing new memories acts as noise
on previously stored ones, and that memory lifetime grows with the square of the number of
synaptic states only when potentiation and depression are precisely balanced. Our synapses
are bounded ($4$-bit counters, clipped to a symmetric band, seven effective states) and the
learning rule applies potentiation and depression in explicitly paired phases. The measured
depression-to-potentiation event ratio is $0.56$--$0.79$ under the configuration used
throughout, and $0.87$--$1.36$ at the higher depression setting. This places the system in
the bounded-synapse regime their analysis addresses; we do not claim it satisfies the
precise potentiation--depression balance their quadratic lifetime result requires.

Their analysis concerns \emph{sequential} storage, in which each new memory perturbs those
already stored, whereas we train all patterns \emph{interleaved to equilibrium} over $40$
epochs. A reader may reasonably object that overwriting is avoided by construction in our
protocol, so finding the substrate intact says nothing about their result. We answer in
three parts. First, we concede the protocol difference and do not claim a palimpsest
result. Second, the comparison we draw is with the prediction that bounded synapses under
continued storage degrade as capacity is approached; we measure the substrate across the
range where behaviour fails and find it intact, which is a statement about our system that
does not require protocol matching. Third, the shared-member result of
Section~\ref{sec:subgroups} responds directly to the objection: if interleaving merely
avoided overwriting by construction, the neurons shared across many assemblies are
precisely where it would show, and they are the group whose density \emph{rises} while the
shared fraction grows from $4.8\pct$ to $64.4\pct$. The substrate is not immune --- density
does fall to $74.27\pct$ and $64.94\pct$ at $90$ and $140$ patterns --- but the behavioural
failure precedes it.

\subsection{Limitations}

This is a measurement study of one architecture on one machine, and one task family (cued
recall of randomly generated column patterns). Whether the dissociation between substrate
and retrieval holds under sequential storage, other learning rules, or other readouts is
untested. The sustained-capacity boundary is soft, with overlapping per-load distributions,
and should not be read as a sharp transition. Throughput figures are reported for context
only; no result above depends on a timing measurement.

% =====================================================================
\section{Methods}
\label{sec:methods}

\subsection{Implementation}

The network is implemented in Objective-C++ and Metal Shading Language and runs on the
integrated GPU of an Apple M1 Pro ($16$-core GPU, $16$\,GB unified memory). Shaders are
compiled at runtime from embedded source, which makes geometry specialisation by string
substitution free; this is used to vary population size without recompiling the host
program.

The configuration used throughout unless stated: $4096$ columns of $32$ neurons; $512$
input columns; $32$ branches per neuron of which $4$ are afferent; $256$ presynaptic
neurons per branch; initial functional synapse density $2^{-4}$; $64$ active columns;
$k=24$ winners per column; dendritic threshold $1$; supralinear branch mode; $64$ stimulus
columns per pattern; contrastive learning with depression applied on alternate free steps;
intrinsic excitability increment $73$ and decrement $1$; presentation length $T=10$ steps
of which the first $5$ are clamped; $40$ training epochs. Seeds vary weight initialisation,
connectivity and the stimulus set jointly. Throughput is $2554$--$2956$ timesteps per
second trained and $3795$--$4930$ untrained, reproducible to $\pm 3\pct$ on clean single
runs.

\subsection{Verification of primitives}

Correctness of the numerical primitives was established before any result was collected.
The SIMD ballot primitive was verified lane-exact at a SIMD width of $32$ against a known
bit pattern. The exact $k$-winners-take-all returns exactly $k$ set bits for
$k \in \{1,3,5,8,16\}$ across all $4096$ columns, at $2.5$\,ns per column. The bit-serial
saturating increment, decrement (borrow path) and symmetric clip were verified against a
scalar reference over $2{,}097{,}152$ random values each, for three clip bands, with zero
mismatches. A parallelised top-$K$ column selection was verified bit-identical to the
serial implementation over $983{,}040$ winner-set words on an untrained network and
$491{,}520$ on a trained one before being adopted.

\subsection{Memory system and sparsity geometry}

Random gather in which each lane follows its own pointer reaches $3$--$5\pct$ of peak
bandwidth, and making the gathered blocks contiguous barely helps. The same random blocks
read cooperatively by all $32$ lanes of a SIMD group reach $102$\,GB/s, a $13.6\times$
improvement at identical block size and identical randomness. All synaptic access is
therefore SIMD-group-cooperative, in tiles of $32$ postsynaptic by $128$ presynaptic binary
weights.

With neuron count, fan-in and spike count held identical, activity clustered into columns
runs at $14{,}585$ timesteps per second against $1{,}285$ for randomly scattered activity
--- an $11.4\times$ difference arising purely from the geometry of the sparse code, because
a presynaptic group of $128$ neurons is silent far more often under columnar activity than
under scattered activity at the same overall rate. Dense binary propagation reaches
$1.36\times10^{12}$ synapse operations per second.

Two alternative execution strategies were measured and rejected. A persistent kernel with a
device-wide barrier is not viable: only $8$ threadgroups are reliably co-resident regardless
of threadgroup size; ordinary device-memory writes are not visible across threadgroups
within a dispatch, so all cross-threadgroup traffic would have to be atomic; and a correct
barrier, which requires the spin read to be a read-modify-write, costs $1.7\,\mu$s against
$2.5\,\mu$s for simply issuing a separate dispatch per timestep. Distributing computation
across the CPU's matrix coprocessor, neural engine and efficiency cores is also ruled out:
a host round trip costs $165\,\mu$s against timesteps of $0.35$--$0.44$\,ms, so any per-step
offload loses.

Per-kernel timing attribution used cumulative ablation in the production encoder. The
device exposes one counter set and does not support counter sampling at dispatch
boundaries, and a marginal-cost method is invalid here because repeated back-to-back
dispatches of the same kernel run from cache.

\subsection{Data and code availability}

The complete implementation (a single Objective-C++ source file containing all Metal
kernels), the benchmark and sweep scripts, and the measurement logs are available at
\url{https://github.com/rakib-nyc/axon}, archived alongside this preprint at
\href{https://doi.org/\doilink}{\doilink}. Every table in this paper is reproducible from a
single command line with the configuration and seed recorded in the corresponding script;
the repository README gives the command for each table.

% =====================================================================
\footnotesize
\begin{thebibliography}{99}
\raggedright

\bibitem{fusiabbott2007}
S.~Fusi and L.~F.~Abbott.
\newblock Limits on the memory storage capacity of bounded synapses.
\newblock \emph{Nature Neuroscience}, 10:485--493, 2007.
\newblock \url{https://doi.org/10.1038/nn1859}

\bibitem{clark2026}
D.~G.~Clark.
\newblock Transient dynamics of associative memory models.
\newblock \emph{Physical Review E}, 113:054301, 2026. arXiv:2506.05303.
\newblock \url{https://arxiv.org/abs/2506.05303}

\bibitem{tamamori2026}
A.~Tamamori.
\newblock Geometric and dynamical analysis of attractor boundaries and storage limits in
kernel Hopfield networks.
\newblock arXiv:2605.00366, 2026.
\newblock \url{https://arxiv.org/abs/2605.00366}

\bibitem{tamamori2025}
A.~Tamamori.
\newblock Quantitative attractor analysis of high-capacity kernel Hopfield networks.
\newblock arXiv:2505.01218, 2025.
\newblock \url{https://arxiv.org/abs/2505.01218}

\bibitem{knoblauch2010}
A.~Knoblauch, G.~Palm, and F.~T.~Sommer.
\newblock Memory capacities for synaptic and structural plasticity.
\newblock \emph{Neural Computation}, 22(2):289--341, 2010.

\bibitem{papadimitriou2020}
C.~H.~Papadimitriou, S.~S.~Vempala, D.~Mitropolsky, M.~Collins, and W.~Maass.
\newblock Brain computation by assemblies of neurons.
\newblock \emph{Proceedings of the National Academy of Sciences}, 117(25):14464--14472,
2020.

\bibitem{makhzani2015}
A.~Makhzani and B.~Frey.
\newblock Winner-take-all autoencoders.
\newblock arXiv:1409.2752, 2015.
\newblock \url{https://arxiv.org/abs/1409.2752}

\bibitem{sacouto2023}
L.~Sacouto and A.~Wichert.
\newblock Competitive learning to generate sparse representations for associative memory.
\newblock \emph{Neural Networks}, 168:32--43, 2023. arXiv:2301.02196.

\end{thebibliography}

\end{document}
